\scanpage{163}
\mspage{78}
\margin{\deleted{\unsure{Dans tous} les topos \deleted{\unsure{envisagés|définis}}\inserted{\unsure{complétés}} sont des cas particuliers des $V_{d_*,d'_*}$ et $V^*_{d_*,d'_*}$ de (260).}}
Je vais faire la liste complète de ces topos\inserted{typiques}, dans le cas le plus élémentaire de tous :
\[
(262)\qquad X=\mathbb E^{1\,\mathrm{an}},\qquad \Theta=\{0\}
\]
on aura donc (en n'écrivant que les $D_{d_*}$ \unsure{str.|et} non vides) les ``drapeaux'' $\{0\},\{1\},\{1,0\}$ et les inclusions de drapeaux non triviales $\{0\}\unsure{-|\longrightarrow}\{0,1\}$, $\{1\}\unsure{-|\longrightarrow}\{0,1\}$.
On aura
\deleted{
\[
\begin{gathered}
D_0=\deleted{\unsure{V_0|D_0}}=\mathbb C
\qquad
\uline{D_0^*=\mathbb C^*}\inserted{=V_0^*\deleted{\unsure{D_0^*}}}
\qquad (V_0^*)
\end{gathered}
\]
\diagram{\begin{tikzcd}[ampersand replacement=\&,column sep=large]
D_1=D_1^* \& \deleted{\unsure{D_1|D_2}}D_{1,0}=D_{1,0}^*=\uline{\text{origine }0\text{ de }\mathbb C}\arrow[l,"\sim"']
\end{tikzcd}}
\[
V_1^{\deleted{\unsure{*}}}=V_{0,\mathbb C}
\]
}
\margin{\deleted{Pour \unsure{éviter les} redondances telles que $D_{d,0}\xrightarrow{\sim}D_d$, on peut convenir que tous les $d_*=(d_0,\ldots,d_r)$ terminent par $d_r=0$, \unsure{mais qu'on n'écrit pas} le dernier $0$; \unsure{pour les opérations} $\frac12$ simp. $D_{d_*}\unsure{-|\longrightarrow}D_{d'_*}$.}}
\[
\left\{
\begin{array}{ll}
D_0=\mathbb C & D_1=\boxed{D_1^*=\inserted{\{P\}}}\deleted{\text{origine }0\text{ de }\mathbb C}\quad\inserted{(\text{origine de }\mathbb C)}\\
\boxed{D_0^*=\mathbb C^*} & {}
\end{array}
\right.
\]
\diagram{\begin{tikzcd}[ampersand replacement=\&,column sep=large,row sep=large]
D_{1,0}\arrow[r,"\sim"]\arrow[d,"{\text{inclusion de }\{0\}\text{ dans }\mathbb C}"'] \& D_1=\{P\} \& D_{1,0}^*=D_{1,0}\\
D_0=\mathbb C \& \&
\end{tikzcd}}
\[
(263)\quad\left\{
\begin{array}{ll}
\boxed{V_0^*\xrightarrow{\sim}V_0\xrightarrow{\sim}D_0^*=\mathbb C^*}
& \text{fibré trivial sur }\mathbb C^*\ (\text{fibre ponctuelle})\\[1ex]
\boxed{V_1=\mathbb D}
& (\deleted{\text{disque}}\inserted{\text{germe}}\text{ de disques ouverts centré en }0)\\[1ex]
\boxed{V_1^*=\mathbb D^*}
& (\text{germe de disques ouverts épointés autour de }0)
\end{array}
\right.
\]
\[
V_{1,0}\xrightarrow{\sim}V_1,\quad V_{1,0}^*\simeq V_1^*
\]
En définitive, il y a ``quatre'' pièces de
\margin{162}
